\documentclass[11pt,a4paper]{article}
\usepackage[margin=25mm]{geometry}
\usepackage[T1]{fontenc}
\usepackage{lmodern}
\usepackage{amsmath,amssymb}
\usepackage{booktabs,longtable,array}
\usepackage{graphicx}
\usepackage{hyperref}
\hypersetup{colorlinks=true,linkcolor=blue,urlcolor=blue}
\setlength{\parskip}{0.5em}
\setlength{\parindent}{0pt}
\setlength{\emergencystretch}{3em}
\newcommand{\NAIRU}{u^{*}}
\newcommand{\gap}{g}
\newcommand{\eps}{\varepsilon}

\title{Japanese SMOG Model\\Results and Coefficient Analysis}
\author{Model output review}
\date{16 September 2026\\Saved model results through 2026Q1}

\begin{document}
\maketitle

\section{Results and charts}
The historical maximum-likelihood sample is 1995Q1--2022Q4 (112 quarters).
The fixed-parameter update extends the state estimates through 2026Q1. At
that date the saved smoothed output gap is $-0.836\%$ of potential output and
the fitted NAIRU is $2.706\%$. Observed unemployment is $2.673\%$, about
0.033 percentage points \emph{below} the NAIRU. The negative output gap
therefore does not align with a simple current-quarter unemployment-minus-
NAIRU reading. The model also contains two lagged unemployment-gap terms, an
inflation signal, and a labour-cost signal, so a one-variable interpretation
would be misleading.

Smoothed states use the full updated sample. Filtered states use observations
available through each plotted quarter while holding the historical
coefficients fixed. The plots are model estimates, not direct measurements
of potential output or the NAIRU.

\begin{figure}[p]
\centering
\includegraphics[width=0.85\linewidth,height=0.28\textheight,keepaspectratio]{\detokenize{../results/figures/jp_smog_output_gap.png}}
\caption{Output gap: smoothed and filtered estimates.}
\end{figure}

\begin{figure}[p]
\centering
\includegraphics[width=0.85\linewidth,height=0.28\textheight,keepaspectratio]{\detokenize{../results/figures/jp_smog_unemployment_and_nairu.png}}
\caption{Observed unemployment and smoothed and filtered NAIRU.}
\end{figure}

\begin{figure}[p]
\centering
\includegraphics[width=0.85\linewidth,height=0.28\textheight,keepaspectratio]{\detokenize{../results/figures/jp_smog_unemployment_gap.png}}
\caption{NAIRU minus observed unemployment: smoothed and filtered estimates.}
\end{figure}

\clearpage
\section{Reading the current output gap}
The model reports a $-0.836\%$ output gap in 2026Q1 even though observed
unemployment is slightly below its fitted NAIRU. Its two lagged
unemployment-gap terms and other signals allow this combination. The latest
local data show core inflation of $1.786\%$ and an expectation input of
$2.600\%$. The smoothed gap contributes $-0.149$ percentage points to the
fitted inflation equation. Its full conditional fitted value is $2.560\%$,
so actual inflation is $0.774$ points lower. This is a decomposition using
the smoothed state, not an out-of-sample forecast.

The negative gap deserves scrutiny, particularly because unemployment is
below the fitted NAIRU and the labour-cost equation imposes a large positive
loading on unemployment slack. But the observed inflation below the
expectation input does not, by itself, show that the gap is too negative or
that inflation pressure is being missed. The inflation residual, the unusual
labour-cost restriction, and the nearly exact unemployment observation
equation all matter for this assessment.

\section{State-space structure}
Japan's eight states are the log output gap $\gap_t$ and its two lags,
potential log output $y_t^*$ and its lag, and the NAIRU $\NAIRU_t$ and its
two lags. The reported output gap is $100\gap_t$, approximately a percentage
deviation from potential. The transition equations are
\begin{align}
 \gap_t &= \phi_1\gap_{t-1}+\phi_2\gap_{t-2}+\eta^g_t,\\
 y_t^* &= y_{t-1}^*+g^*+\eta^y_t,\\
 \NAIRU_t &= \NAIRU_{t-1}+\eta^u_t.
\end{align}
The gap is AR(2), potential output has a quarterly log drift $g^*$, and the
NAIRU is a random walk. The code uses independent Gaussian observation and
state shocks with diagonal covariance matrices. Boosted HP trends from the
historical data set the known initial state, with zero initial covariance.
They also help construct optimizer starting values; Python estimates the
coefficients from the Japanese signals.

\section{Observation equations}
The four measured signals are log GDP $y_t$, unemployment $u_t$, core
inflation $\pi_t$, and annual labour-cost growth $n_t$. Inflation
expectations $\pi_t^e$ come from the Tankan-labelled input. The import-price
term $m_{t-1}$ is a lagged four-quarter level change. Let $D_t=1$ from
2013Q1 onward and zero before that. The implemented equations are
\begin{align}
 y_t &= y_t^*+\gap_t+\eps^y_t,\\
 u_t &= \NAIRU_t+\lambda_3\gap_t
       +\rho_1(u_{t-1}-\NAIRU_{t-1})
       +\rho_2(u_{t-2}-\NAIRU_{t-2})+\eps^u_t,\\
 \pi_t &= D_t\{(1-\beta_1-\beta_2)\pi_t^e
                 +\beta_1\pi_{t-1}+\beta_2\pi_{t-2}\}\notag\\
 &\quad +(1-D_t)\{(1-\beta_1-\beta_2-\beta_3-\gamma)\pi_t^e
          +\beta_1\pi_{t-1}+\beta_2\pi_{t-2}
          +\beta_3\pi_{t-3}+\gamma n_{t-1}\}\notag\\
 &\quad +\lambda_1\gap_t+\psi m_{t-1}+\eps^\pi_t,\\
 n_t &= (1-\beta_4)\pi_t^e+\beta_4\pi_{t-1}
       +L(u_{t-1}-\NAIRU_{t-1})+\eps^n_t,\\
 L &= 8+\frac{1}{1+\exp(-\lambda_2)}.
\end{align}
The labour-cost loading is constrained to lie strictly between 8 and 9.
It is $L$, rather than the raw coefficient $\lambda_2$, that multiplies the
unemployment gap. A one percentage point output gap corresponds to
$\gap_t=0.01$ in the inflation and unemployment equations.

\section{Fitted coefficients and interpretation}
\begin{longtable}{@{}>{$}l<{$}r p{0.65\linewidth}@{}}
\caption{Japanese maximum-likelihood coefficients. Values are rounded.}\\
\toprule
\text{Parameter} & \multicolumn{1}{c}{\text{Estimate}} & Meaning in the implemented model\\
\midrule
\endfirsthead
\toprule
\text{Parameter} & \multicolumn{1}{c}{\text{Estimate}} & Meaning\\
\midrule
\endhead
\beta_1 & 0.770747 & Weight on one-quarter lagged inflation in both regimes.\\
\beta_2 & -0.091067 & Weight on two-quarter lagged inflation in both regimes.\\
\beta_3 & -0.044532 & Weight on three-quarter lagged inflation before 2013Q1.\\
\beta_4 & 2.176609 & Weight on lagged inflation in labour-cost growth; the implied expectation weight is $-1.1766$.\\
\gamma & -0.010815 & Loading of lagged labour-cost growth in pre-2013 inflation.\\
\lambda_1 & 17.859318 & Inflation loading on the log output gap; a one percentage point output gap adds about 0.1786 percentage points through this term.\\
\lambda_2 & 6.884643 & Raw parameter transformed into the positive labour-gap loading $L=8.99898$; it is not the direct loading itself.\\
\lambda_3 & -4.011514 & Unemployment loading on the log output gap; a one percentage point gap contributes about $-0.0401$ percentage points.\\
\rho_1 & 1.304940 & Loading of the first lagged unemployment-minus-NAIRU gap.\\
\rho_2 & -0.482033 & Loading of the second lagged unemployment-minus-NAIRU gap.\\
\phi_1 & 1.210914 & First AR coefficient of the latent output gap.\\
\phi_2 & -0.319172 & Second AR coefficient of the latent output gap.\\
\psi & 0.002892 & Inflation loading of the lagged four-quarter import-price change.\\
g^* & 0.001678 & Quarterly potential-log-output drift, about 0.67\% at an annualized log-growth rate.\\
\eps_1 & 0.006487 & Standard deviation of log GDP measurement noise, about 0.649\% in log-output terms.\\
\eps_2 & 0.000010 & Standard deviation of unemployment measurement noise, in percentage points; at the imposed lower bound.\\
\eps_3 & 0.285608 & Standard deviation of inflation measurement noise, in percentage points.\\
\eps_4 & 5.349199 & Standard deviation of labour-cost-growth measurement noise, in percentage points.\\
\eps_5 & 0.008176 & Standard deviation of the log output-gap state shock, about 0.818\% of potential output.\\
\eps_6 & 0.005007 & Standard deviation of the potential-log-output state shock.\\
\eps_7 & 0.112964 & Standard deviation of the quarterly NAIRU state shock, in unemployment percentage points.\\
\bottomrule
\end{longtable}

The implied inflation-expectation weights are $0.3203$ from 2013Q1 and
$0.3757$ before 2013. The labour equation's $L\approx8.999$ is a large,
positive loading: holding its other terms fixed, one percentage point more
unemployment above the previous NAIRU raises predicted labour-cost growth
by about nine percentage points. That sign is unusual for a slack measure.
It is inherited from the Japanese specification and constrained to the
8--9 interval, so its magnitude should be read as a structural modelling
choice, not an unrestricted empirical estimate. The labour-cost signal also
has high measurement noise. The near-zero $\eps_2$ makes unemployment almost
an exact measurement conditional on the latent states and is another reason
to treat the NAIRU and output gap as specification-sensitive.

\section{Likelihood fit and convergence}
SciPy's bounded L-BFGS-B optimizer maximizes the Kalman-filter log
likelihood. The accepted-iteration CSV records all 21 parameters and their
full-sample likelihood; a second CSV records each quarter's contribution.
The Japanese fit converged after 441 accepted iterations:
\begin{center}
\begin{tabular}{@{}lr@{}}
\toprule
Fit statistic & Japan\\
\midrule
Historical quarters & 112\\
Initial log likelihood & $-137.5153$\\
Iteration 25 & 23.8740\\
Iteration 50 & 24.9822\\
Final log likelihood & 25.7183\\
\bottomrule
\end{tabular}
\end{center}
The NAIRU innovation scale, $\eps_7=0.112964$, is above the imposed
data-derived floor of $0.072897$. The unemployment observation noise,
$\eps_2=0.000010$, is at its lower bound. Convergence establishes a numerical
optimum from the selected starting point and bounds; it does not establish
a global optimum, tight identification, or out-of-sample forecast accuracy.
The script does not export standard errors or forecast scores. Quarter-level
likelihood terms are Kalman one-step-ahead predictive densities, not direct
squared-error measures; early missing regressors cause signals to be masked.

\section{Reproduction and source files}
From this model folder, run the generation script to re-estimate historical
coefficients and overwrite the history CSVs. Then run the update script to
apply the fixed fit to later quarters. This report uses the generation code,
fitted parameter JSON, both fit-history CSVs, the updated state CSV, the
\texttt{Update} tab of the local input workbook, and the saved figures.
Its numerical tables are a snapshot and do not change automatically when
Python is rerun.

\end{document}
